% =====================================================================
% REMAINING CHANGES — checked against VeRedact-2.tex (3,406 lines)
%
% DONE already (nothing to do): Table I, notation row X_i^red,
% cost-notation, computation-cost, communication-cost (+ sentence),
% primitives, gas tables; Evaluation intro (B1); paragraph after the
% computation-cost table (B2).
%
% LEFT: 14 edits below, in file order. Each says LINES (in VeRedact-2.tex),
% the OLD text exactly as it is there, and the NEW text to paste.
% Line numbers shift after an edit, so work from the BOTTOM up
% (edit 14 first, edit 1 last) and they stay correct.
% =====================================================================


% ---------------------------------------------------------------------
% EDIT 1 — LINE 365 (Related Work) — double full stop
% OLD:  cross-chain accountability~\cite{ref25,ref28,ref29,ref32,ref35}.. Verifiable
% NEW:
cross-chain accountability~\cite{ref25,ref28,ref29,ref32,ref35}. Verifiable


% ---------------------------------------------------------------------
% EDIT 2 — LINE 674 (Phase 1, Step 1, equation for PP) — COMPILE ERROR
% OLD:  \end{aligned}}
% NEW:  (one closing brace removed)
\end{aligned}


% ---------------------------------------------------------------------
% EDIT 3 — LINE 817 (Phase 2, Step 1, equation for D_i) — purple "2" still printed
% OLD:  \rho_i\stackrel{\$}{\leftarrow}\{0,1\}^{\textcolor{purple}{2}\lambda}}
% NEW:  choose ONE
% (a) drop the 2, consistent with K_idx and K_A at lines 709-710:
\rho_i\stackrel{\$}{\leftarrow}\{0,1\}^{\lambda}}
% (b) keep the 2 (salt longer than the keys) but in black:
% \rho_i\stackrel{\$}{\leftarrow}\{0,1\}^{2\lambda}}


% ---------------------------------------------------------------------
% EDIT 4 — LINES 2739-2740 (end of Computation Cost Analysis)
% OLD:  below quantify these effects using identical PQ primitives across
%       schemes.
% NEW:
below quantify these effects with every baseline at its own classical
instantiation at comparable (about 128-bit) security, so the measured
gap includes the full cost of post-quantum protection in VeRedact-PQ.


% ---------------------------------------------------------------------
% EDIT 5 — LINES 2871-2880 (Experimental Setup, 3rd paragraph)
% OLD:  the whole paragraph from
%       "To separate architectural effects from primitive choice, all baselines"
%       to "classical instantiations are reported for reference in Experiment~4."
% NEW:
All baselines were reimplemented from their papers within the same
framework, on the same dataset, request trace, and ledger, and follow
their original workflows, including per-request authorization and
adaptation. Because none of the baseline papers defines a post-quantum
variant of its primitives (RSA accumulators and tags, pairing-based
chameleon hashes, or attribute-based encryption), each baseline is
instantiated with its own classical primitives at approximately 128-bit
security: BLS12-381 pairings for Schemes~\cite{ref1} and~\cite{ref34},
3072-bit RSA for Schemes~\cite{ref1}, \cite{ref13}, and~\cite{ref34},
and secp256k1 for Schemes~\cite{ref1} and~\cite{ref27}. A stage that a
baseline's paper does not define is reported as not supported rather
than synthesized.


% ---------------------------------------------------------------------
% EDIT 6 — LINES 2947-2948 (Experiment 1, first paragraph)
% OLD:  is compared with Refs.~\cite{ref14}, \cite{ref17}, \cite{ref27},
%       and~\cite{ref33}. In addition, three internal variants isolate the
% NEW:
is compared with Schemes~\cite{ref1}, \cite{ref13}, \cite{ref27},
and~\cite{ref34}. In addition, three internal variants isolate the


% ---------------------------------------------------------------------
% EDIT 7 — LINES 2977-2981 (Experiment 1, first results paragraph)
% OLD:  the arrival rate until each scheme saturates. Refs.~\cite{ref14}
%       and~\cite{ref17} saturate early because every request requires an
%       independent chameleon-hash adaptation and checkpoint update, while
%       Refs.~\cite{ref27} and~\cite{ref33} additionally incur per-request
%       threshold or policy-based authorization. VeRedact-PQ sustains the
% NEW:
the arrival rate until each scheme saturates. All four baselines
execute one chameleon-hash adaptation and one ledger write per request:
Scheme~\cite{ref13} additionally regenerates block tags, updates its
accumulator, and rebuilds the block's Merkle tree, Schemes~\cite{ref1}
and~\cite{ref27} run a $t$-party adaptation per request, and
Scheme~\cite{ref34} decrypts its policy ciphertext per request.
Scheme~\cite{ref1} further rejects every request that targets an already
redacted block, so its finalized redactions are bounded by the number of
distinct target blocks. VeRedact-PQ sustains the


% ---------------------------------------------------------------------
% EDIT 8 — LINES 3017-3019 (Experiment 2, first paragraph)
% OLD:  verification. VeRedact-PQ is compared with Liu \emph{et al.}~\cite{ref27},
%       Dong \emph{et al.}~\cite{ref15}, and Li \emph{et al.}~\cite{ref1}, which
%       support distributed or threshold redaction authorization. An internal
% NEW:
verification. VeRedact-PQ is compared with Scheme~\cite{ref1},
whose $t$ approving full nodes each validate the request, and
Scheme~\cite{ref34}, whose redactor decrypts a policy ciphertext with
attribute keys issued by distinct verification nodes; the committee
size $n$ is mapped to the number of full nodes in Scheme~\cite{ref1} and
to the number of policy attributes in Scheme~\cite{ref34}.
Scheme~\cite{ref13} (a single trusted system manager) and
Scheme~\cite{ref27} (threshold enforced only inside the adaptation)
define no separate authorization step and are therefore not included.
An internal


% ---------------------------------------------------------------------
% EDIT 9 — LINE 3038 (Experiment 2, results (a))
% OLD:  baselines perform committee interaction for every request, so their
% NEW:
baselines authorize every request individually, so their


% ---------------------------------------------------------------------
% EDIT 10 — LINES 3046-3047 (Experiment 2, results (b))
% OLD:  increases with committee size for all schemes because more PQ
%       signatures must be generated and verified. For VeRedact-PQ, this growth
% NEW:
increases with committee size for all schemes because more signatures,
or more attribute pairings in Scheme~\cite{ref34}, must be processed.
For VeRedact-PQ, this growth


% ---------------------------------------------------------------------
% EDIT 11 — LINES 3061-3063 (Experiment 3, first paragraph)
% OLD:  control evidence sharing. VeRedact-PQ is compared with
%       VRBC~\cite{ref13}, Xue \emph{et al.}~\cite{ref33}, and
%       Miao \emph{et al.}~\cite{ref20}. An internal variant,
% NEW:
control evidence sharing. VeRedact-PQ is compared with
Scheme~\cite{ref13}, whose audit returns aggregated tags and an
aggregated accumulator non-membership witness for the challenged
blocks, and Scheme~\cite{ref1}, whose consistency check returns each
queried block with an accumulator witness; Schemes~\cite{ref27}
and~\cite{ref34} define no audit protocol. An internal variant,


% ---------------------------------------------------------------------
% EDIT 12 — LINES 3108-3110 (Experiment 4, first paragraph)
% OLD:  verification. VeRedact-PQ is compared with Refs.~\cite{ref13},
%       \cite{ref20}, and~\cite{ref33} in both PQ-adapted and original
%       classical instantiations. To evaluate verification granularity,
% NEW:
verification. VeRedact-PQ is compared with Schemes~\cite{ref13}
and~\cite{ref1} in their original classical instantiations. To evaluate
verification granularity,


% ---------------------------------------------------------------------
% EDIT 13 — LINES 3135-3145 (Experiment 4, results (b)) — whole paragraph
% OLD:  from "Fig.~\ref{fig:exp4_verify}(b) shows the verification-time breakdown."
%       to   "checks of Phase~6."
% NEW:
Fig.~\ref{fig:exp4_verify}(b) shows the verification-time breakdown.
Compared with the classical baselines, VeRedact-PQ incurs higher
absolute signature-verification cost, which quantifies the overhead of
post-quantum protection. Nevertheless, VeRedact-PQ reduces the number
of PQ verifications by verifying shared authorization evidence once per
batch. Because each returned record carries individual membership and
execution evidence, every injected modified, substituted, or stale
record is rejected individually while all valid records remain
acceptable, as stale evidence fails the version, epoch, and checkpoint
checks of Phase~6. Scheme~\cite{ref1} also decides per record, whereas
the aggregated proof of Scheme~\cite{ref13} yields one decision for all
challenged blocks, so a single injected record causes the valid records
of the same response to be rejected as well.


% ---------------------------------------------------------------------
% EDIT 14 — LINES 3154-3155 (Experiment 5, first paragraph) — matches the new gas-table row
% OLD:  policy and committee registration (Phase~1), batch-checkpoint anchoring
%       $A_b$ (Phase~2), redaction finalization, which atomically commits the
% NEW:
policy and committee registration (Phase~1), batch-checkpoint anchoring
$A_b$ (Phase~2), authorization anchoring $H(Auth_e^B)$ (Phase~4),
redaction finalization, which atomically commits the
